\section{Introduction}
\label{sec:introduction}

The class $\QAC$ consists of constant-depth quantum circuits with
arbitrary one-qubit gates and generalized Toffoli gates of unbounded
arity, using polynomially many total qubits. Ancillary qubits are initialized
to zero.
The supports of gates in a layer are pairwise disjoint. This last
condition makes the availability of fanout a substantive question:
one cannot simultaneously use a qubit as the control of arbitrarily
many two-qubit gates. Moore asked whether the coherent fanout
operation can be implemented in this model, and conjectured a
negative answer. Its unitary parity counterpart is equivalent
by Hadamard conjugation
\citep{Moore1999,GreenHomerMoorePollett2002}.

We prove Moore's parity conjecture in the measured-output model and
therefore also exclude coherent fanout. The lower bound
holds for bounded-error computation with a single measured output,
even when the computation changes the input and leaves arbitrary
garbage in the other registers.

\subsection{The circuit model and the result}

For $x\in\{0,1\}^{n}$, let
\[
  \Parity_n(x)=x_1\mathbin{\oplus}\cdots\mathbin{\oplus}x_n.
\]
An allowed circuit acts on $\lvert x\rangle\lvert0\rangle^{\otimes a}$.
Its gates are arbitrary one-qubit unitaries and generalized Toffoli
gates
\[
 \lvert z_1,\ldots,z_t,b\rangle
 \longmapsto
 \lvert z_1,\ldots,z_t,b\mathbin{\oplus}(z_1\cdots z_t)\rangle,
 \qquad t\ge1,
\]
on distinct control and target qubits. Each gate has unit cost,
regardless of its arity. A depth-$d$ circuit is a product of $d$
layers of gates with pairwise disjoint supports. There is no
geometric restriction. At the end, one designated qubit is measured
in the computational basis. The other registers may be discarded.
Intermediate measurements, postselection, and additional fanout
gates are not allowed. All gates and their complex entries may
depend arbitrarily on $n$, but not on $x$.

\begin{theorem}
\label{thm:parity}
Fix positive integers $d,k,C$ and a real $0<\varepsilon\le1/2$.
There is an integer $n_0=n_0(d,k,C,\varepsilon)$ such that,
for every $n\ge n_0$, every allowed circuit of depth at most $d$
on at most $C(n+1)^k$ total qubits has an input $x\in\{0,1\}^{n}$ on which its answer equals
$\Parity_n(x)$ with probability strictly less than $1/2+\varepsilon$.
\end{theorem}

Here total qubits means the $n$ inputs together with all ancillary
qubits. Disjoint nonempty gate supports imply that a depth-$d$ circuit
on $N$ qubits contains at most $dN$ gates. Thus the theorem covers the
usual fixed-depth model with polynomially many gates and polynomially
many ancillas: the ancillary bound gives a polynomial bound on $N$,
and a polynomial bound on $N$ already bounds the number of gates.

In particular, no nonuniform constant-depth family with polynomial
resources computes parity with any fixed positive worst-case advantage.
The theorem is eventual for each fixed depth, polynomial bound, and advantage;
it permits arbitrary finite exceptions at small input lengths.

To make the connection with fanout explicit, let
\[
 F_n\lvert y_1,\ldots,y_n,b\rangle
 =\lvert y_1\mathbin{\oplus}b,\ldots,
          y_n\mathbin{\oplus}b,b\rangle.
\]
For the Hadamard gate $H$, conjugating a controlled-NOT by
$H\otimes H$ interchanges its control and target; this follows
directly from its action on the two-qubit basis. Writing $F_n$
as the product of its commuting controlled-NOT operations gives
\[
 H^{\otimes(n+1)}F_nH^{\otimes(n+1)}
 \lvert y,b\rangle
 =\lvert y,b\mathbin{\oplus}\Parity_n(y)\rangle.
\]
Thus a constant-depth polynomial-resource implementation of
$F_n$, together with two layers of one-qubit gates and a zero
target, would contradict Theorem~\ref{thm:parity}. This also
excludes Moore's original formulation, in which the ancillary
qubits must be restored to zero \citep{Moore1999}.

Theorem~\ref{thm:parity} also has state-preparation consequences
corresponding to those of Gretta, Gupta, and Joshi
\citep[Corollary~4.10 and Theorem~4.14]{GGJ2026Fourier}.
Let $\rho$ be the retained $n$-qubit state of an allowed unitary
preparation from all-zero inputs, with arbitrary final ancillas discarded.
Its felinity is $2\sum_y p_y p_{\bar y}$, where
$p_y=\langle y|\rho|y\rangle$ and $\bar y$ is the bitwise complement of $y$.
At each fixed depth and fixed polynomial total-qubit bound, for every fixed
$A>0$, eventually every such preparation has felinity less than $n^{-A}$.
For every fixed $0<\delta<1$, eventually every such preparation and every
integer $k$ with $n^\delta\le k\le n/2$ also satisfy
\[
 \tfrac12\bigl\|\rho-|D_k^n\rangle\langle D_k^n|\bigr\|_1
 >\frac1{80k},
 \qquad
 |D_k^n\rangle=\binom nk^{-1/2}\sum_{|y|=k}|y\rangle.
\]
Here trace distance is one half of the trace norm. The companion states
the same bounds
\citep[Corollary~8 (State-preparation obstructions), Section~7.1]{OpenAILocalization2026}.
The argument after the parity proof in Section~\ref{sec:parity}
justifies the pointwise reductions, including exact promotion with
arbitrary final garbage and polynomial total-qubit overhead.

\subsection{Previous work}

The classical class $\mathsf{AC}^{0}$ allows unbounded-fan-in AND and
OR gates, NOT gates, and unrestricted reuse of wires. Parity was among
the first explicit functions shown to require superpolynomial size
at constant depth, by Furst, Saxe, and Sipser and independently Ajtai
\citep{FurstSaxeSipser1984,Ajtai1983}; H\aa stad obtained almost optimal
exponential bounds through the switching lemma \citep{Hastad1986}.
The quantum question combines the additional expressive power of
one-qubit unitaries with the disjoint-support restriction on each layer.
Moore's formulation and the counting classes studied by Green, Homer,
Moore, and Pollett made the role of coherent fanout explicit
\citep{Moore1999,GreenHomerMoorePollett2002}. Constant-depth
constructions with fanout further demonstrated the computational
importance of that operation \citep{HoyerSpalek2005}.

Early quantum lower bounds exposed the importance of the ancillary
register. Fang, Fenner, Green, Homer, and Zhang proved a depth bound of
order $\log(n/(a+1))$ for exact parity computation with $a$ ancillas
that must be restored to zero
\citep[Corollary~4.6]{FangFennerGreenHomerZhang2006}.
Pad\'e, Fenner, Grier, and Thierauf ruled out exact clean simulation
at entangling depth two \citep{PadeFennerGrierThierauf2020}; Fenner,
Grier, Pad\'e, and Thierauf subsequently removed the final-state
cleanliness requirement for exact single-output parity at that depth
\citep[Theorem~6.1]{FennerGrierPadeThierauf2025}.
Here and in the depth-specific comparisons below, \emph{entangling
depth} counts only layers of multiqubit gates; arbitrary one-qubit
layers may occur between them. This convention differs from our total
layer count by at most a constant factor for fixed depth.

Rosenthal obtained lower bounds for approximate parity and fanout
unitaries and related cat-like states, together with constant-depth
approximation constructions using superpolynomial resources
\citep{Rosenthal2021}. Approximation of a parity unitary is a stronger
task than computing one measured classical bit: the latter may destroy
the input and leave arbitrary garbage. Nadimpalli, Parham, Vasconcelos,
and Yuen developed Pauli-spectrum methods with depth-dependent
restrictions on the ancillary count
\citep{NadimpalliParhamVasconcelosYuen2024}. Anshu, Dong, Ou, and Yao
obtained lower bounds allowing barely superlinear ancillary counts
\citep{AnshuDongOuYao2025}; Dong, Ou, and Yao developed further
low-degree estimates and applications to quantum channels
\citep{DongOuYao2025}. These resource restrictions do not cover an
arbitrary fixed polynomial ancillary bound.

Joshi, Tal, Vasconcelos, and Wright ruled out exact parity at
entangling depth three for sufficiently large input lengths, with no
bound on the number of gates or ancillas. They also proved approximate
parity lower bounds at entangling depth two for one measured output,
without requiring input preservation
\citep{JoshiTalVasconcelosWright2025}. Kintali's September 2026
preprint reports exponentially small average parity correlation at
entangling depth at most three, independently of the number of gates
and ancillas, and an exponential gate lower bound for exact parity at
entangling depth at most four
\citep[Theorems~1--2]{Kintali2026ShallowParity}.\footnote{Joshi et al.
acknowledge ChatGPT's contribution to a lemma's proof idea and other
interactions; Kintali discloses assistance from Claude, ChatGPT, and
two open-weight models. Neither source specifies model versions.}
These depth-specific statements do not cover every fixed depth.
Gretta, Gupta, and Joshi characterize the general parity question
through Fourier concentration \citep{GGJ2026Fourier}, while Xu and Li
extend related fanout reductions to symmetric functions
\citep{XuLi2026Fanout}. Our theorem permits every fixed depth and
every polynomial total-qubit bound simultaneously. All structural
statements needed for its proof are established below.

The companion article \emph{Product-projection localization and the
QAC0 parity lower bound} proves a uniform operator-norm localization
theorem and an independent parity lower bound
\citep[Sections~2 and~6]{OpenAILocalization2026}. The present article instead
proves state-dependent propagation, projection insertion, and pruning
for regular trajectories. Both arguments use product-projection counts
and elementary Chebyshev filters; neither structural theorem is used
to prove the other.

\subsection{The proof and its main tools}

A generalized Toffoli gate is a reflection $I-2P$, where $P$ is a
product of one-qubit rank-one projections. This representation suggests
counting how many local factors of a product test fail. A single such
gate can couple arbitrarily many qubits, so its conjugation need not
preserve small operator support. Moreover, product-projection filters
can have large norms away from the count range where they approximate
well. Controlling the state on which each approximation is used is
therefore essential.
Informally, a trajectory sequence is \emph{regular} when, for every
sufficiently small collection of reflections in a layer, the combined count
of failures of their local factors has a rapidly decaying tail on the state
entering that layer. The precise definition is asymptotic and allows
arbitrary gates and states at each circuit size.

Polynomial locality and concentration have been related in work of
Kuwahara, Arad, Amico, and Vedral and of Anshu and Metger
\citep{KuwaharaAradAmicoVedral2017,AnshuMetger2023}. In the circuit
setting, Pauli concentration and low-degree approximation are central
to the works just discussed. An operator has degree at most $\ell$ when it is a sum of operators
supported on at most $\ell$ sites each. We retain the elementary link
between this degree and the number of product-basis coordinates an
operator can change. The additional steps here are two-sided tail transport and
stability under a product projection, which permit a trajectory to be
made regular by a state-dependent pruning rule.

The first tool, Proposition~\ref{prop:propagation}, transports
operators of low degree along a regular
trajectory, with an arbitrarily small fixed loss in the degree
exponent. Its conclusion is an approximation on the trajectory
vector. The second tool, Proposition~\ref{prop:projection-insertion},
shows that regularity survives insertion
of an arbitrary product projection followed by reverse evolution.
The key step is Lemma~\ref{lem:tail-transport}, which compares two
regular trajectories. One polynomial detects a tail at the starting end;
another approximates the product projection at the finishing end.
They are applied on different sides of an inner product, so their
large global norms are paid for by the appropriate tail estimates.

Projection stability then makes regularity available without an
initial assumption on the circuit. Starting from a product state,
we retain a reflection when its projection has probability at least
$m^{-6}$ on the current state, where $m$ is the number of qubits.
The retained layer is regular, and the deleted gates change the
state by at most $2m^{-2}$ per layer. Iterating this rule yields a
regular trajectory close to the original evolution
(Proposition~\ref{prop:pruning}). The stability
and propagation statements concern arbitrary product tests and
regular trajectories, rather than parity alone.

Finally, we copy each basis input into a reference register that the
subsequent computation leaves untouched.
On the uniform coherent input, worst-case parity success becomes
a signed acceptance expectation bounded away from zero. The
pruned circuit has almost the same expectation. Write $W$ for that circuit,
$\psi_0$ for the initial product state, and $P_B$ for the accepting output
projection. The copying gates survive pruning, so conjugating the reference
sign back through $W$ sends $\psi_0$
to a vector with exactly $n$ excitations relative to the initial product state
$\psi_0$. The signed expectation pairs this vector with $W^\dagger P_BW\psi_0$,
whose tail for that same count is bounded by
Corollary~\ref{cor:product-tail}. Because the total number $m$ of qubits is polynomially
bounded in $n$, a sufficiently small fixed $h>0$ can be chosen
with $m^h=o(n)$, separating these two scales and giving the contradiction.

Sections~\ref{sec:framework} and~\ref{sec:filters} introduce the
formalism and the polynomial filters. Section~\ref{sec:propagation}
proves low-degree propagation, and Section~\ref{sec:transport}
establishes tail transport. Projection stability and pruning are
proved in Section~\ref{sec:pruning}; Section~\ref{sec:parity}
completes the parity argument.
